
There are three boundaries in that description.  The existing coding agents are
clients of the harness: they author and review work, but they do not carry the
control model in their prompts.  Around them, the harness maintains the belief
state, constructs and scores the candidate actions, imposes the act-gates, and
records why a choice was made.  Outside both sits the shared substrate --- the
files, commits, job events, and typed graph records against which a claim can be
checked.  A claim of success closes the loop only when the substrate is changed
and re-observed through a path that the author of the claim does not control.

Conceptually, one cycle has six stages.  \emph{Perceive} turns heterogeneous
artifacts into a structured observation.  \emph{Believe} updates operational
hypotheses and the trust placed in each evidence channel.  \emph{Evaluate}
predicts the consequences of candidate moves and keeps the active-inference
score separate from engineering augmentations.  \emph{Select} constructs a
bounded action space, may abstain, and records which term governed the choice.
\emph{Act} is not a single model utterance: it includes construction, authoring,
bounded repair, independent review, a grounded write, and re-observation.
\emph{Learn} uses the independently witnessed outcomes to recalibrate later
choices and, between runs, to revise the structures from which future
candidates are proposed.  The six names describe the control cycle; the
timestamped phases in Section~\ref{sec:catalog} show how that cycle expands in a
real execution.

Figure~\ref{fig:aif-brain} is a control map rather than a deployment diagram.
Its five columns give the main cycle room to read from left to right: perceive,
believe, evaluate, select, act, and re-observe.  The band below shows the
assurance and learning apparatus that makes a decision independently checkable
rather than a well-organised opinion.  Prior preferences (the model's C
vector) are not a drawn node and carry no pattern number.

The figure is generated at build time from the same wiring registries the
harness is checked against, and its footer pins the registry revisions it was
drawn from.  Its edge classes separate three kinds of claim that a
hand-maintained diagram would conflate: wiring that is \emph{drawn} (including
wiring drawn but carrying no data, and wiring retired by a recorded decision),
routes that are \emph{measured} in recorded runs and anchored in the trace
store, and edges the \emph{theory} expects, classed by whether the code
realises them.  Agreement between the drawing and the running system is
therefore a checked property of the build, not an illustration.  The
R-numbers align directly with the catalogue in Section~\ref{sec:catalog}.

R10's scheduled entrypoint has an unattended, reviewed, grounded witness on
record and was subsequently operator-disabled; that records an operator
choice, not a lower maturity grade.  Table~\ref{tab:code-audit} records how
Campaign~S closed the code-audit items (R11, R12, R15, R17).
These are maturity claims about the named paths, not claims that their models
are generally effective.

\begin{center}
\begin{minipage}{\linewidth}
\captionsetup{font=footnotesize,skip=3pt}
\captionof{table}{Closure of the code-audit items by Campaign S.}
\label{tab:code-audit}
\centering
\scriptsize
\setlength{\tabcolsep}{3pt}
\renewcommand{\arraystretch}{0.96}
\begin{tabular}{@{}lp{2.4cm}p{10.9cm}@{}}
\toprule
\textbf{R-\#} & \textbf{Status} & \textbf{Live evidence} \\
\midrule
R12 & Run-evidenced & Both separately recorded layers passed on both ticks; independent review and grounded re-observation supplied Layer~2. \\
R11 & Run-evidenced & Both ticks selected one item from each live field at feasible total cost 7; replay reproduced both portfolios. \\
R17 & Run-evidenced & The inter-tick pass reduced a frozen 26-observation model in 20~ms, replayed exactly, and fed Tick B. \\
R15 & Run-evidenced & Tick A's witnessed success changed Beta(1,1) to Beta(2,1) and exploration to consolidation before Tick B. \\
\bottomrule
\end{tabular}
\end{minipage}
\end{center}

\hypertarget{pat-r19}{}
\begin{figure}[h]
\centering
\includegraphics[width=\linewidth]{aif-control-map-live}
\caption{Checked topology of the active-inference harness, generated at build
time from the wiring registries and pinned to their source revisions in the
footer.  Nodes carry the catalogue's R-numbers; R3a is the prediction-error projection
that measured routes pass through between structured observation (R2) and
evidence-channel precision (R7) --- no direct R2-to-R7 call exists, so the
measured path is R2$\rightarrow$R3a$\rightarrow$R7.  Edges are wiring claims of
stated standing: solid edges are drawn wiring (transition, constraint, or
unclassified), faint dotted edges are drawn but carry no data, and dash-dot
edges are retired by a recorded decision; purple, teal, and orange edges are
added by decision --- found in code, intended by a ruling, or conditional on a
guard; amber, crimson, and olive edges are theory classes from the conformance
registry --- realised in code but undrawn, not realised, and path-dependent;
green edges were observed in recorded routes, anchored in the trace store.}
\label{fig:aif-brain}
\end{figure}
\FloatBarrier

The map also has a model-side incompleteness layer that a topology drawing
alone cannot show.  In particular, a generic Lean carrier or a formula that
accepts a caller-supplied predictive distribution does not establish that the
running machine constructs that distribution from its beliefs, model, and
policy.  Table~\ref{tab:q-interface-gaps} is regenerated from the checked
interface registry on every paper build, keeping the missing constructor,
underpowered deliveries, and next wiring actions visible alongside the map.

\input{sec-q-interface-generated}
